% Original analytic teaching diagram, 2026-11-09.
\definecolor{VB}{HTML}{7188A3}\definecolor{VBF}{HTML}{E5EAF0}
\definecolor{VG}{HTML}{829C8B}\definecolor{VGF}{HTML}{E8EEE9}
\definecolor{VY}{HTML}{B4A16F}\definecolor{VYF}{HTML}{F1ECDD}
\definecolor{VR}{HTML}{A7787D}\definecolor{VRF}{HTML}{F0E5E6}
\definecolor{VP}{HTML}{9588A5}\definecolor{VPF}{HTML}{ECE8F0}
\definecolor{VI}{HTML}{3B4252}\definecolor{VM}{HTML}{6B7280}\definecolor{VL}{HTML}{4C4D4F}
\begin{tikzpicture}[x=1mm,y=1mm,>=Stealth,font=\figurefont\fontsize{8}{10}\selectfont,
lab/.style={text=VM,font=\figurefont\fontsize{7.5}{9}\selectfont},
box/.style={rounded corners=1.2mm,line width=.8pt}]
\draw[box,draw=VB,fill=VPF] (0,0) rectangle (84,46);
\draw[box,draw=VB,fill=VYF] (5,4) rectangle (70,38);
\draw[box,draw=VB,fill=VGF] (10,8) rectangle (56,30);
\draw[box,draw=VB,fill=VBF] (15,12) rectangle (42,22);
\node[text=VI] at (28.5,17) {线性规划};
\node[text=VI,anchor=south] at (33,23.2) {凸二次规划};
\node[text=VI,anchor=south] at (37.5,31.2) {二阶锥规划};
\node[text=VI,anchor=south] at (42,39.2) {半正定规划};
\node[lab,anchor=west] at (43,17) {$\bm{Q}=\bm{0}$};
\node[lab,anchor=west] at (57,25) {上图形式};
\node[lab,anchor=west] at (71,33) {Schur补};
\node[draw=VR,fill=VRF,box,inner sep=1.6mm,text=VI] (n) at (102,17) {非凸二次问题};
\draw[-{Stealth[length=2mm]},VL,line width=.8pt,densely dashed] (n.north) |- (84,44);
\node[lab,anchor=west] at (86.5,46.2) {半正定松弛};
\end{tikzpicture}
